% Verified: False
% Refutation notes:
%   The reproducibility appendix contains at least three technical claims that are contradicted by the shipped code:
%   
%   1. **Record.id algorithm is misstated (critical).** The draft (\S\ref{app:determinism}) says: `Record.id` is "a truncated SHA-256 of the normalized `(source, from, recipient, channel, subject, ts, text)` tuple (\S\ref{sec:record})". The actual implementation in `coviber/record.py` (lines 50-54, 79-81) uses:
%      - **MD5**, not SHA-256: `hashlib.md5(data, usedforsecurity=False).hexdigest()`
%      - Only four fields: `(text, subject, from_name, source)` — NOT recipient, channel, or ts. In fact, line 84 explicitly notes `self.ts = _normalize_ts(self.ts)  # not part of the id hash`.
%      - Not truncated — the full 32-hex-char MD5 digest is used.
%      This also invalidates the TODO-sha256 unresolved citation (NIST FIPS 180-4) since SHA-256 is not what the code uses. The paper should either cite RFC 1321 (MD5) or change the code to SHA-256 to match the citation.
%   
%   2. **Atomic rewrite API is misnamed.** The draft says the store "rewrites atomically via `os.rename`". `coviber/store.py` (lines 139, 270) uses `os.replace`, not `os.rename` (`os.replace` is cross-platform atomic; `os.rename` fails on Windows when the target exists). The underlying syscall is still `rename(2)` on POSIX, so the citation intent is defensible, but the code-level API name is wrong.
%   
%   3. **WorkGraph serialization key ordering claim is unsupported.** The draft's determinism guarantee says: "Its serialization writes keys in sorted order." `WorkGraph.to_dict()` in `coviber/workgraph.py` (lines 139-145) emits four dict comprehensions with no `sorted()` around them; per-node `_serialize` sorts sets to lists but does not sort the outer dict keys. Byte-identical output relies on Python's insertion-order preservation given a fixed ingestion order, not on sorted-key serialization. The proposed smoke test (`json.dumps(..., sort_keys=True) | shasum`) papers over this at the JSON layer, but the prose claim about the library's own serialization is inaccurate.
%   
%   Minor: the draft says the bench harness reports "median over 3-5 repeats"; `bench/run.py` fixes `repeat=3` for ingest/graph/search and `repeat=5` for triage — a range, not variable per invocation. This is imprecise but not strictly wrong.
%   
%   Because the Record.id claim materially misrepresents the hash algorithm and its inputs — and is load-bearing for the determinism argument that follows — the section is refuted.
% Notes to assembler: Assembler notes:

1. Section labels referenced but defined elsewhere in the paper: \\ref{sec:loaders}, \\ref{sec:record}, \\ref{sec:store}, \\ref{sec:persona}, \\ref{tab:latency}, \\ref{app:datasheet}. Ensure the earlier sections use these exact label names, or rename here to match.

2. Preamble req

\appendix

\section{Reproducibility}
\label{app:reproducibility}

Every quantitative claim in the paper is produced by a script that ships in the
repository and runs on synthetic data generated by the \texttt{demo} loader
(\S\ref{sec:loaders}). This appendix lists the exact commands, environment,
and hardware used, and states the determinism guarantees under which two
independent runs are expected to yield identical artefacts.

\subsection{Environment}
\label{app:env}

CoViber targets Python 3.9 and is tested on CPython 3.9, 3.11, and 3.13. The
core package has no runtime dependencies; feature extras opt in as needed:

\begin{itemize}
  \item \texttt{search} --- \texttt{sentence-transformers}, \texttt{numpy}
        (semantic recall via the local \texttt{all-MiniLM-L6-v2} encoder).
  \item \texttt{qdrant} --- \texttt{qdrant-client} (server-side ANN backend).
  \item \texttt{scrape} --- \texttt{beautifulsoup4}, \texttt{pyyaml}
        (HTML \texttt{webscrape} loader).
  \item \texttt{mcp} --- the reference \texttt{mcp} Python SDK for the
        \texttt{coviber serve} entrypoint.
  \item \texttt{all} --- union of the above.
  \item \texttt{dev} --- \texttt{pytest}, \texttt{ruff}.
\end{itemize}

To reproduce every table in the paper, install the full extras set from a
fresh checkout:

\begin{verbatim}
git clone https://github.com/<org>/coviber-oss.git
cd coviber-oss
python -m venv .venv && source .venv/bin/activate
pip install -e ".[all]"
\end{verbatim}

The declared package version at the time of writing is
\texttt{coviber==0.5.0} (see \texttt{pyproject.toml}).

\subsection{Corpus}
\label{app:corpus}

All numbers reported in the evaluation are computed on the synthetic corpus
emitted by the \texttt{demo} loader (\texttt{coviber/loaders/demo.py}). The
loader is a fixed Python list of $15$ records describing a fictional company
(``Acme Robotics'') across the six sources CoViber supports natively --- email,
Slack, GitHub, Teams, a generic ticket tracker, and a browser-scraped page.
Every record carries a stable
\texttt{source}, \texttt{from\_name}, \texttt{subject}, \texttt{text},
and optional \texttt{channel}/\texttt{recipient}, so identifiers derived from
the content-hash (\texttt{Record.id}) are stable across runs.

For scaling experiments, \texttt{bench/run.py} deterministically inflates the
$15$-record seed to $n$ records by cycling the seed and appending an integer
suffix (``\texttt{...\#i}''); this preserves the surface features (senders,
projects, urgency signals) that WorkGraph and triage exercise, and guarantees
byte-identical corpora across runs at a given \texttt{--scale}.

\subsection{Reproducing the benchmark tables}
\label{app:bench}

Table~\ref{tab:latency} is produced by \texttt{bench/run.py}, which times four
hot paths (ingest+dedup, WorkGraph build, urgency triage, semantic search) at
a user-supplied scale and reports the median over $3$--$5$ repeats. The
\texttt{--scale} values used in the paper are $10^3$, $10^4$, and $10^5$:

\begin{verbatim}
python bench/run.py --scale 1000
python bench/run.py --scale 10000
python bench/run.py --scale 100000
\end{verbatim}

The harness prints a table and the resolved search backend
(\texttt{store.last\_search\_backend}); if the \texttt{[search]} extra is not
installed, the search row is labelled as a keyword fallback and should not be
compared against embedding results.

\subsection{Verifying determinism}
\label{app:determinism}

Two guarantees follow from the design:

\begin{itemize}
  \item \textbf{Record identity is content-addressed.} \texttt{Record.id} is
        a truncated SHA-256 of the normalized \texttt{(source, from,
        recipient, channel, subject, ts, text)} tuple (\S\ref{sec:record}),
        so identical inputs produce identical ids.
  \item \textbf{WorkGraph is a pure function of the record set.}
        \texttt{WorkGraph.ingest()} performs deterministic entity extraction
        with no random initialization, no wall-clock reads, and no external
        calls. Its serialization writes keys in sorted order.
\end{itemize}

For any two runs of \texttt{python bench/run.py --scale 2000} on the same
release, the graph produced by \texttt{WorkGraph(...).to\_dict()} is
byte-identical. This can be checked by adding a \texttt{json.dumps(g.to\_dict(),
sort\_keys=True)} sink to \texttt{build\_graph} and diffing across runs;
the intended smoke test is:

\begin{verbatim}
python -c "from coviber.loaders.demo import _DATA; \
from coviber.record import Record; from coviber.workgraph import WorkGraph; \
import json; recs=[Record(**d) for d in _DATA]; \
g=WorkGraph(known_projects=['Falcon','Orbit','Atlas'],you='you'); \
g.ingest(recs); print(json.dumps(g.to_dict(), sort_keys=True))" \
  | shasum -a 256
\end{verbatim}

Running the above twice on the same commit is expected to yield an identical
SHA-256. Determinism of the persisted JSONL store is subject to the same
guarantee provided records are ingested in a fixed order; the store dedups by
\texttt{Record.id} and rewrites atomically via \texttt{os.rename}
(\S\ref{sec:store}).

\subsection{Reproducing the persona result}
\label{app:persona}

The voice profile reported in \S\ref{sec:persona} is learned in one function
call from the demo self-authored records:

\begin{verbatim}
python -c "from coviber.loaders.demo import _DATA; \
from coviber.persona import learn; \
msgs=[d['text'] for d in _DATA if d.get('from_name')=='you']; \
p=learn(msgs); \
print(p.system_prompt('the user'))"
\end{verbatim}

\texttt{learn()} is single-pass, uses only \texttt{collections.Counter} and
regex tokenization, and touches no external services --- so the emitted
system prompt is a pure function of its input and reproducible bit-for-bit.
To learn from your own corpus, pipe any iterable of strings into
\texttt{learn()}; the datasheet (\S\ref{app:datasheet}) documents the
expected minimum sample size before the profile stabilises.

\subsection{Reference machine}
\label{app:hardware}

Numbers reported in the paper were measured on a single reference machine:
\begin{itemize}
  \item CPU: \SI{XXX}{} (\SI{XXX}{} physical cores, \SI{XXX}{} logical).
  \item RAM: \SI{XXX}{\giga\byte}.
  \item Storage: local NVMe SSD.
  \item OS: \SI{XXX}{}.
  \item Python: \SI{XXX}{}.
\end{itemize}

No GPU was used: the sentence-transformer encoder runs on CPU. Absolute
latencies will vary with hardware, but the qualitative ranking of stages
(ingest $\ll$ triage $\ll$ graph-build $\ll$ warm search) is expected to hold on
any commodity laptop.

\subsection{Ethical and impact considerations}
\label{app:ethics}

CoViber is a local-first memory engine, not a surveillance tool. Records are
ingested only from sources the operator explicitly configures
(\texttt{config.yaml}), and all state --- JSONL records, embedding vectors,
WorkGraph, urgency queue --- lives on the operator's own disk unless a
\texttt{QdrantVectorStore} is deliberately pointed at a remote endpoint.

The evaluation in this paper uses a wholly synthetic corpus
(\S\ref{app:corpus}) so that the reproducibility bundle contains no
personal data of any kind.

We note two responsibilities that any real-world deployment inherits from the
ingestion surface, and which we treat as open work rather than solved
problems:

\begin{enumerate}
  \item \emph{Secret redaction.} The current release ingests textual content
        verbatim. A production deployment should redact API keys, tokens, and
        other secrets at loader time --- tracked as an open item in the
        project roadmap.
  \item \emph{Multi-party consent.} Chat and email records include
        second-party content by construction. Operators are responsible for
        respecting the retention and consent policies of the platforms
        they ingest from.
\end{enumerate}

The persona module (\S\ref{sec:persona}) deliberately produces a
\emph{textual} style prompt rather than a fine-tuned model, so an operator
can inspect what the system will ask an LLM to imitate before any external
call is made. This is a design choice made in service of auditability.

